\documentclass[10pt,a4paper]{amsart}
\usepackage[margin=19mm]{geometry}
\usepackage{amsmath,amssymb,mathtools}
\usepackage[scr=rsfs,cal=euler]{mathalfa}
\usepackage{xfrac}
\usepackage[unicode,hidelinks,bookmarks=false]{hyperref}
\newcommand{\ie}{i.e.\ }
\newcommand{\cf}{cf.\ }
\newcommand{\resp}{respectively\ }
\newcommand{\Subsection}{\subsection}
\providecommand{\nobreakdash}{\nobreak\hbox{-}}
\setlength{\emergencystretch}{2em}
\multlinegap0pt
\usepackage{amssymb}
\usepackage{mathtools}
\usepackage{dsfont}
\usepackage{float}
\usepackage{enumerate}
\usepackage{algorithm}\usepackage{algpseudocode}
\newtheorem{ass}{Assumption}
\newtheorem{lemma}{Lemma}
\newcommand{\E}{\mathbb{E}}
\newcommand{\N}{\mathbb{N}}
\newcommand{\R}{\mathbb{R}}
\title{Preconditioning for the high-order sampling of the invariant distribution of parabolic semilinear SPDEs}
\author{Charles-Edouard Br\'ehier, Adrien Busnot Laurent, Arnaud Debussche, Gilles Vilmart}
\date{Source: arXiv:2512.17714v1 (CC BY 4.0)}
\begin{document}
\maketitle

\noindent\emph{Source and licence.} Preconditioning for the high-order sampling of the invariant distribution of parabolic semilinear SPDEs. Version arXiv:2512.17714v1 (CC BY 4.0), distributed by arXiv under the Creative Commons Attribution 4.0 International licence, http://creativecommons.org/licenses/by/4.0/, as recorded in the arXiv licence metadata for this exact version and shown on the abstract page https://arxiv.org/abs/2512.17714v1. Full licence notice: \texttt{licenses/arxiv-2512.17714-CC-BY-4.0.txt}.\par\smallskip

\noindent\textit{Editorial scope.} Counted unit: the article's lemma lem:order1 (source0.tex lines 1177--1189). The article's complete original proof of it is delivered with it (source0.tex lines 1191--1205, one \texttt{proof} environment of 15 lines). No mathematics is written, repaired or re-derived: the statement and the proof are the article's own lines, in the article's own notation, and no retained line is substituted or re-flowed.  Retained uncounted as the source-local context the proof needs: lines 150--152; lines 430--436; lines 1115--1135; lines 1138--1144.  Nothing was omitted from any retained span, so the package carries no omission record.  No retained line contains a citation, so the wrapper carries no bibliography and no bibliography entry.  No retained line contains a figure environment, an image-inclusion command or drawing code, so no image is delivered and the package declares no asset dependency.  Editorial formatting only, in addition: an AMS \texttt{amsart} preamble that restates the article's own declarations and macros used by the retained lines, namely the environments \texttt{ass}, \texttt{lemma} on separate counters, together with the standard packages those lines need.  The retained environments print in this document as Ass 1, Lemma 1 in order of appearance, whereas the article prints the same items with the numbering of its own full document.  The complete native source of the article as distributed in the arXiv e-print for this exact version is bundled unchanged as \texttt{sources/191339/source0.tex}.
\begin{equation} \label{eq:def_mu_star}
d\mu_\star(x)=Z^{-1}\exp\bigl(-2V(x)\bigr)d\nu(x),
\end{equation}

\begin{align}
\label{eq:precondSPDE}
dY(t)&=\big(-Y(t)+QF(Y(t))\big)dt+dW^{Q}(t),\\
\label{eq:precondSPDE_K}
dY^K(t)&=\big(-Y^K(t)+QF^K(Y^K(t))\big)dt+dW^{Q,K}(t),\quad
Y^K(0)=\pi^K Y(0).
\end{align}

The aim of this section is to create, under stronger regularity assumptions, high order methods for sampling $\mu_\star$, given by \eqref{eq:def_mu_star}.
We consider one step integrators with postprocessors,
\[
Y_{n+1}^{\Delta t}=\psi^{\Delta t}(Y_n^{\Delta t}),
\quad
\overline{Y}_n^{\Delta t}=\overline{\psi}^{\Delta t}(Y_n^{\Delta t}).
\]
%
We consider a vector field $F\in \mathfrak{X}^{4,\beta}_P$ and test functions $\varphi\in\mathcal{C}^{6,\beta}_P$.
%
Recall that the dynamics $\bigl(Y(t)\bigr)_{t\ge 0}$ to be discretized is given by the preconditioned SPDE~\eqref{eq:precondSPDE}. This process is ergodic and mixing, and its unique invariant distribution is $\mu_\star$.
Thus, we assume a similar property for the integrator.
We consider the following assumptions.
\begin{ass}
\label{ass:ergodic_num}
For some $\Delta t_0>0$, if $\Delta t\in(0,\Delta t_0]$, the Markov process $\bigl(Y_n^{\Delta t}\bigr)_{n\ge 0}$ is ergodic and mixing. Its unique invariant distribution, denoted by $\mu^{\Delta t}$, satisfies
\[
\E[\varphi(Y_n)]\underset{n\to\infty}\to \int \varphi d\mu^{\Delta t},
\]
for all bounded and continuous test functions $\varphi:H\to \R$, and any arbitrary initial condition $y\in H$.
\end{ass}

\begin{ass}
\label{ass:bounded_mom}
The integrator has bounded moments of all order: for any $\kappa\in\N$, there exists $C>0$ such that
\[
\underset{\Delta t\in(0,\Delta t_0]}\sup~\underset{n\ge 0}\sup~\E[|Y_n^{\Delta t}|_\beta^\kappa]\le C(1+|Y_0|_\beta^\kappa).
\]
\end{ass}

\begin{lemma}\label{lem:order1}
Under Assumptions \ref{ass:ergodic_num} and \ref{ass:bounded_mom}, assume that, for any $\varphi\in\mathcal{C}^{4,\beta}_P$, one has the Taylor-type expansion
\[
\E\bigl[\varphi\bigl(\psi^{\Delta t}(y)\bigr)\bigr]-\varphi(y)=\Delta t\mathcal{L}\varphi(y)+\Delta t^2R_1(\varphi,y,\Delta t),
\]
for all $\Delta t$ assumed small enough,
with $\underset{\Delta t\in(0,\Delta t_0]}\sup~|R_1(\varphi,y,\Delta t)|\le C(1+|y|_\beta^\kappa)$ for some $p\in\N$ and $C>0$.
%
Then, for $\Delta t \rightarrow 0$ and $\varphi\in\mathcal{C}^{4,\beta}_P$,
\begin{equation}\label{eq:expansion_order1}
\int\varphi d\mu^{\Delta t}-\int\varphi d\mu_\star={\rm O}(\Delta t).
\end{equation}
\end{lemma}

\begin{proof}[Proof of Lemma~\ref{lem:order1}]
Since $\mu^{\Delta t}$ is the invariant distribution of the Markov chain $\bigl(Y_n^{\Delta t}\bigr)_{n\ge 0}$, one has the following equality: for all $\Delta t\in(0,\Delta t_0]$,
\begin{equation}\label{eq:equilibrium}
\int \Bigl(\E\bigl[\varphi\bigl(\psi^{\Delta t}(y)\bigr)\bigr]-\varphi(y)\Bigr)d\mu^{\Delta t}(y)=0.
\end{equation}
Then, one obtains
\[
0=\int \mathcal{L}\varphi(y)d\mu^{\Delta t}(y)+\Delta t\int R_1(\varphi,y,\Delta t) d\mu^{\Delta t}(y).
\]
Choosing $\varphi=\Psi$ then yields
\[
\int\varphi d\mu^{\Delta t}-\int\varphi d\mu_\star=\Delta t\int R_1(\varphi,y,\Delta t)d\mu^{\Delta t}(y)={\rm O}(\Delta t),
\]
owing to the moment bounds for the numerical invariant distribution $\mu^{\Delta t}$.
\end{proof}

\end{document}
